\chapter{Kleisli Iteration and Delayed Feedback}\label{ch:kleisli}

\paragraph{Step as a Kleisli morphism.}
A morphism $M = [S, \mathrm{step}] : A \to B$ of Mol is a Mealy machine
with state space $S$. The function
$\mathrm{step} : S \times A \to B \times S$ is, under the curry
isomorphism, a Kleisli morphism of the state monad $T_S$ on $S$
(\autoref{thm:5}, \autoref{sec:prelim-monads}). The response of $M$ on
input words is the iterate of that Kleisli morphism
(\autoref{thm:33}). Sequential composition in Mol is the product of
state spaces (\autoref{def:comp}); it coincides with Kleisli
composition of $T_S$ precisely when one factor is pure.

\paragraph{Classes.}
A molecule is \emph{pure} when its state space is isomorphic to $()$
(\autoref{def:puremol}), \emph{effectful over} $\mathrm{Ctx}$ when its
state space is isomorphic to $\mathrm{Ctx}$ (\autoref{def:effmol}), and
\emph{hybrid} otherwise (\autoref{def:hybridmol}). Hybrid is the
residual class: Mol composition of two effectful molecules has state
space $\mathrm{Ctx} \times \mathrm{Ctx}$ and is therefore hybrid
(\autoref{thm:67}).

\paragraph{Delayed feedback.}
The reactor loop $\mathrm{Tr}^{U}_{A,B}$ of \autoref{def:reactor} feeds
the $U$-output of a body back into its $U$-input as a register. The
construction satisfies vanishing, superposing, and sliding of pure maps,
proved from the machine definition as \autoref{thm:34}--\autoref{thm:36}.
It fails the Joyal--Street--Verity yanking axiom: $\mathrm{Tr} (\sigma)$
is the one-step delay on $U$, not the identity. Mol therefore carries a
\emph{delayed trace}, not a JSV traced monoidal structure. Hoisting
(\autoref{thm:37}) uses sliding of pure maps to shrink loop-carried
state.

\setcatalog{33}
\begin{theorem}[Kleisli iteration of step (\autoref{thm:33})]\label{thm:33}
\textbf{[C].} Let $M = [S, \mathrm{step}] : A \to B$ be a morphism of
Mol, and let $T_S$ be the state monad $T_S X = (X \times S)^{S}$ of
\autoref{sec:prelim-monads}.
\begin{enumerate}
\item[(i)] Under the curry isomorphism,
$\mathrm{step}^{\sharp} : A \to T_S B$ is a Kleisli morphism of $T_S$.
\item[(ii)] The response $r_{M,s} : A^{*} \to B^{*}$ of
\autoref{def:response} is the unique stream map obtained by iterating
that Kleisli morphism from $s$: $r_{M,s} (\varepsilon) = \varepsilon$
and
\[
r_{M,s} (a \cdot w) = b \cdot r_{M,s'} (w)
\quad\text{where}\quad (b, s') = \mathrm{step} (s, a).
\]
\item[(iii)] Conversely, every Kleisli morphism $A \to T_S B$ determines
a unique Mol class with representative $(S, \mathrm{step})$.
\item[(iv)] If $\varphi : S \to S'$ is a state-space isomorphism of two
representatives of the same Mol morphism, the corresponding Kleisli
morphisms are conjugate by $\varphi$.
\end{enumerate}
The state space of a Mol morphism is well defined up to bijection by
\autoref{def:mol}; behavioral equivalence (\autoref{def:bisim}) may
identify machines with non-bijective state spaces, and minimization
(\autoref{ch:coalgebra}) produces a generally smaller state space with
the same responses.
\end{theorem}

\begin{proof}
\emph{(i).} A Kleisli morphism $A \to T_S B$ is a function
$A \to (B \times S)^{S}$, which by uncurrying is a function
$S \times A \to B \times S$. That is exactly $\mathrm{step}$.

\emph{(ii).} The clauses are the definition of response
(\autoref{def:response}). Uniqueness of the stream map is induction on
word length: the empty word is fixed, and each letter determines a
unique output letter and a unique successor state because
$\mathrm{step}$ is a total function (\autoref{def:machine}).

\emph{(iii).} Given a Kleisli morphism $k : A \to T_S B$, uncurry to a
step function and take the class $[S, \mathrm{uncurry} (k)]$. Two Kleisli
morphisms that uncurry to state-space isomorphic machines determine the
same class (\autoref{thm:2}); if they uncurry to the same function they
are equal.

\emph{(iv).} The isomorphism hypothesis is
$\mathrm{step}' (\varphi (s), a) = (b, \varphi (s'))$ whenever
$\mathrm{step} (s, a) = (b, s')$. Currying both sides yields
$k' (a) (\varphi (s)) = (b, \varphi (s'))$ whenever $k (a) (s) = (b, s')$,
which is conjugacy of the Kleisli morphisms by $\varphi$.
\end{proof}

\paragraph{Remark (dummy-wire factorization).}
One may write a representative $(S, \mathrm{step})$ as a Mol composite
$q \circ e \circ p$ in which $p : A \to A \otimes S$ is the pure map
$a \mapsto (a, s_0)$, $q : B \otimes S \to B$ is the pure projection, and
$e$ is a machine with state $S$ that ignores its $S$-input wire and
writes the successor state onto its $S$-output wire. The construction
equals $M$ because $e$ ignores the dummy wire. It does not give a
canonical dataflow of the state, and it does not identify $S$ among
behavioral equivalents. The Kleisli reading of (i)--(ii) is the one used
below.

\setcatalog{34}
\begin{theorem}[Sliding of pure maps (\autoref{thm:34})]\label{thm:34}
\textbf{[C].} Let $M : A \otimes U \to B \otimes U$ be a molecule with
state space $S$ and let $f : B \to B'$ and $g : A' \to A$ be pure. Then
\[
\mathrm{Tr}^{U}_{A,B'}\bigl ((f \otimes 1_U) \circ M\bigr)
= f \circ \mathrm{Tr}^{U}_{A,B} (M),
\qquad
\mathrm{Tr}^{U}_{A',B}\bigl (M \circ (g \otimes 1_U)\bigr)
= \mathrm{Tr}^{U}_{A,B} (M) \circ g.
\]
A pure map on the feed-forward wires of a delayed trace may be slid
completely out of the loop.
\end{theorem}

\begin{proof}
Write $\mathit{step}_M (s, (a, u)) = ((b, u'), s')$. The reactor loop
$\mathrm{Tr}^{U}_{A,B} (M)$ has state $S \times U$ and
\[
\mathit{step}'\bigl ((s, u), a\bigr) = \bigl (b, (s', u')\bigr)
\]
(\autoref{def:reactor}).

\emph{Output sliding.} The machine $(f \otimes 1_U) \circ M$ has state
$() \times S \cong S$. On input $(a, u)$ in state $s$ it runs $M$,
obtaining $((b, u'), s')$, then applies the pure map $f$ to $b$,
obtaining $((f (b), u'), s')$. Its reactor loop therefore has state
$S \times U$ and emits $f (b)$ while advancing to $(s', u')$. The
composite $f \circ \mathrm{Tr}^{U}_{A,B} (M)$ has state
$() \times (S \times U) \cong S \times U$ and, on the same input, emits
$f (b)$ after the same register update. The canonical bijection of state
spaces is a state-space isomorphism (\autoref{def:stiso}), so the two
morphisms are equal in Mol.

\emph{Input sliding.} The machine $M \circ (g \otimes 1_U)$ has state
$S \times ()$. On input $(a', u)$ it applies $g$ to $a'$ and then $M$.
The reactor loop emits the same $b$ and stores the same $u'$ as
$\mathrm{Tr} (M)$ run on $g (a')$. The composite
$\mathrm{Tr} (M) \circ g$ does the same, up to the unit-state bijection.
\end{proof}

\setcatalog{35}
\begin{theorem}[Vanishing and delay (\autoref{thm:35})]\label{thm:35}
\textbf{[C].} Let $\mathrm{Tr}$ be the reactor loop of
\autoref{def:reactor}. Then:
\begin{enumerate}
\item[(i)] \emph{Vanishing I.} For $M : A \otimes I \to B \otimes I$,
$\mathrm{Tr}^{I}_{A,B} (M) = M$ in Mol, via $S \times () \cong S$.
\item[(ii)] \emph{Vanishing II.} For
$M : A \otimes (U \otimes V) \to B \otimes (U \otimes V)$,
\[
\mathrm{Tr}^{U \otimes V}_{A,B} (M)
= \mathrm{Tr}^{U}_{A,B}\Bigl (
  \mathrm{Tr}^{V}_{A \otimes U,\, B \otimes U} (M) \Bigr)
\]
in Mol, via $S \times (U \times V) \cong (S \times V) \times U$.
\item[(iii)] \emph{Delay, in place of yanking.} Let
$\sigma_{U,U} : U \otimes U \to U \otimes U$ be the braiding of
\autoref{def:braid}, and let $\Delta_U : U \to U$ be the \emph{delay}
molecule with state space $U$ and
$\mathit{step}_{\Delta} (u, x) = (u, x)$. Then
$\mathrm{Tr}^{U}_{U,U} (\sigma_{U,U}) = \Delta_U$. In particular
$\mathrm{Tr} (\sigma_{U,U})$ equals the identity on $U$ if and only if
$U$ is a singleton, so the reactor loop fails the Joyal--Street--Verity
yanking axiom whenever $|U| > 1$ \citep{joyalsv}.
\end{enumerate}
\end{theorem}

\begin{proof}
\emph{(i).} The reactor loop of $M$ has state $S \times ()$ and, writing
$\mathit{step}_M (s, (a, ())) = ((b, ()), s')$, its transition is
$((s, ()), a) \mapsto (b, (s', ()))$. The bijection
$\varphi (s, ()) = s$ is a state-space isomorphism onto $M$.

\emph{(ii).} Direct trace: $\mathrm{Tr}^{U \otimes V} (M)$ has state
$S \times (U \times V)$. On input $a$ in state $(s, (u, v))$ it runs one
step of $M$ on $(a, (u, v))$, obtaining $((b, (u', v')), s')$, and
advances to $(s', (u', v'))$.

Inner-then-outer: view $M$ as a body
$(A \otimes U) \otimes V \to (B \otimes U) \otimes V$. The inner loop
$\mathrm{Tr}^{V} (M)$ has state $S \times V$; on input $(a, u)$ in state
$(s, v)$ it emits $(b, u')$ and stores $v'$. The outer loop
$\mathrm{Tr}^{U}$ of that machine has state $(S \times V) \times U$; on
input $a$ in state $((s, v), u)$ it emits $b$ and stores
$((s', v'), u')$. The bijection
$\Phi (s, (u, v)) = ((s, v), u)$ carries the direct-trace transition to
the inner-then-outer transition, and both emit $b$. Hence the two
machines are state-space isomorphic.

\emph{(iii).} The braiding is pure:
$\mathit{swap} ((), (x, y)) = ((y, x), ())$. Taking $A = B = U$ and
feeding the second factor back, the reactor loop has state $() \times U
\cong U$. On input $x$ in register $u$,
\[
\mathit{swap} ((), (x, u)) = ((u, x), ()),
\]
so the loop emits $u$ (the old register) and stores $x$ (the current
input). That is the transition of $\Delta_U$. If $|U| > 1$, pick
distinct $u, x \in U$; delay started at $u$ emits $u$ on input $x$,
while $\mathrm{id}_U$ emits $x$. The two responses differ, so
$\Delta_U \neq \mathrm{id}_U$ in Mol.
\end{proof}

\setcatalog{36}
\begin{theorem}[Superposing (\autoref{thm:36})]\label{thm:36}
\textbf{[C].} Let $M : A \otimes U \to B \otimes U$ and let
$N : C \to D$ be any molecule. Then
\[
\mathrm{Tr}^{U}_{A \otimes C,\, B \otimes D} (M \otimes N)
= \mathrm{Tr}^{U}_{A,B} (M) \otimes N
\]
in Mol. A delayed trace running beside an independent molecule is the
loop and the molecule side by side.
\end{theorem}

\begin{proof}
Let $M$ have state $S$ and $N$ have state $T$. The tensor
$M \otimes N : (A \otimes U) \otimes C \to (B \otimes U) \otimes D$
reassociates to a body of type $(A \otimes C) \otimes U \to
(B \otimes D) \otimes U$ with state $S \times T$
(\autoref{def:tensor}). Its reactor loop has state
$(S \times T) \times U$ and, on input $(a, c)$ in state
$((s, t), u)$, runs $M$ on $(a, u)$ and $N$ on $c$ independently,
emitting $(b, d)$ and storing $((s', t'), u')$.

The tensor $\mathrm{Tr} (M) \otimes N$ has state $(S \times U) \times T$.
On input $(a, c)$ in state $((s, u), t)$ it runs the loop of $M$ on $a$
and $N$ on $c$, emitting $(b, d)$ and storing $((s', u'), t')$. The
bijection $\Phi ((s, t), u) = ((s, u), t)$ is a state-space isomorphism
of the two machines.
\end{proof}

\setcatalog{37}
\begin{theorem}[Hoisting and loop-state footprint (\autoref{thm:37})]\label{thm:37}
\textbf{[N].} Let $M : A \otimes U \to B \otimes U$ be the body of a
reactor loop, presented by a machine $(S, \mathrm{step})$, and suppose
the loop state carries purely redundant information: $U = U' \otimes U''$
with a pure map $h : U' \to U''$ such that the register is initialized
in the graph $i = \langle 1_{U'}, h\rangle : U' \to U$ and every register
value in the run lies in that graph, and the loop body preserves the
graph:
$M \circ (1_A \otimes i) = (1_B \otimes i) \circ M'$ for the compressed
body $M' : A \otimes U' \to B \otimes U'$. Then:
\emph{(a) Hoisting.} $\mathrm{Tr}^{U} (M) \simeq \mathrm{Tr}^{U'} (M')$
(\autoref{thm:6}); the pure redundancy $h$ is hoisted out of the loop, and
the loop-state footprint shrinks from $|U| = |U'| \cdot |U''|$ to $|U'|$.
\emph{(b) Monotonicity.} With fixed-size state records of $r$ bytes, the
resident cache footprint of a loop with loop state $W$ is $|W| \cdot r$
bytes, strictly increasing in $|W|$; hence the cache footprint of loop
state is monotone in state size, and the hoisted loop of (a) is strictly
smaller whenever $|U''| > 1$.
\end{theorem}

\begin{proof}
\emph{(a).} The loop $\mathrm{Tr}^{U} (M)$ carries state space $S \otimes
U$; its register is initialized in the graph, and since the body preserves
the graph, the register stays in the graph at every step, so all reachable
states lie in $S \otimes i (U')$. Conjugate by the bijection $\varphi :
S \otimes i (U') \to S \otimes U'$, $\varphi (s, (u', h (u'))) = (s, u')$.
On a reachable state with input $a$ the full loop computes
$M\bigl (a, (u', h (u'))\bigr) = (1_B \otimes i)\bigl (M'(a, u')\bigr)$, so it
emits the same output and advances to $(s', (u'_{1}, h (u'_{1})))$, whose
image under $\varphi$ is $(s', u'_{1})$; the compressed loop
$\mathrm{Tr}^{U'} (M')$ on $(s, u')$ with input $a$ emits the same output
and advances to $(s', u'_{1})$: the register evolves by the loop body's
transition in both loops (\autoref{ch:trace}). The two machines are
therefore behaviorally equivalent, $\mathrm{Tr}^{U} (M) \simeq
\mathrm{Tr}^{U'} (M')$ (\autoref{thm:6}), and their minimal realizations
(\autoref{ch:coalgebra}) are isomorphic. The compressed loop stores only
$u'$; the value $u'' = h (u')$ is recomputed by the pure map $h$ where it
is consumed, outside the loop, so the resident loop state shrinks from
$|U'| \cdot |U''|$ to $|U'|$ elements.

\emph{(b).} Fixed-size records of $r$ bytes make the resident working set
of a loop with loop state $W$ exactly $|W| \cdot r$ bytes, a strictly
increasing function of the state cardinality $|W|$. Replacing $U$ by $U'$
as in (a) multiplies the footprint by the factor $1/|U''| \leq 1$, and
strictly so whenever $|U''| > 1$: the cache footprint of loop state is
monotone in state size.
\end{proof}

\paragraph{Hoisting in practice.}
Together the theorems give the operational shape of a hot loop. By
\autoref{thm:33} the step of a stage is a Kleisli morphism of its own
state monad, iterated on the input stream; pure stages have unit state
and hold none. By \autoref{thm:34} and \autoref{thm:36} a delayed
feedback loop can be rearranged so that pure work and independent side
computations live outside the loop, and by \autoref{thm:37} loop state
that is purely determined by the remainder of the loop state is
compressed away. Yanking is unavailable: closing a swap produces a
delay (\autoref{thm:35}), so a register that is swapped onto a
feed-forward wire is emitted one step late. The loop that runs on the
hot path is the one with the smallest resident state after compression.
This is what the software standard's policies [ALLOC] and [CACHE]
enforce: no allocation in the hot path, the loop state small enough to
stay resident. The following chapters quantify the rest:
\autoref{ch:coalgebra} makes the state minimal among behavioral
equivalents (\autoref{thm:38}), and \autoref{ch:trace} bounds the
number of iterations such a loop needs (\autoref{thm:42}).
